\section{Related Work}
\label{sec:related}

\textbf{Safety \& Alignment.} A large body of work has examined improved methods for alignment~\citep{rafailov2023direct,ethayarajh2024kto,zou2023representation,bai2022constitutional,ouyang2022training,touvron2023llama-2,team2024gemma}. While we tie alignment approaches to potential downstream jailbreaks, we do so mainly be examining aligned artifacts which go through more rigorous alignment procedures than most other open source models. We rely on the Gemma~\citep{team2024gemma} and Llama-2~\citep{touvron2023llama-2} base and aligned models throughout this work, as the safety alignment built in these two models are closest to the technology applied in frontier proprietary models.

\textbf{Jailbreaking Methods.} A large body of work has examined methods for jailbreaking aligned LLMs, including leveraging finetuning, decoding strategies, prefilling strategies, optimization strategies, and even persuasion~\citep{andriushchenko2024jailbreaking,zou2023universal,qi2023fine,huang2023catastrophic,zeng2024johnny,zhan2023removing,yang2023shadow,gade2023badllama,prefilling-attack,qi2023visual}. Many approaches try to handle jailbreaking through a systems approach by monitoring inputs and outputs with machine learning models~\citep{inan2023llama}, but this is only as good as the monitoring mechanism (which can also be jailbroken)~\citep{break-llama-guard}.

\textbf{Superficial Alignment Hypothesis and Per-token Effects of Alignment Fine-tuning.} Our work is closely related to the Superficial Alignment Hypothesis~(SAH)~\cite{zhou2023lima}, which posits that the alignment process for current LLMs only superficially changes the formats of inputs and outputs to be used for interacting with the users. %It is also relevant to \citet{lin2024unlocking}, which also finds that the alignment's effects decay for later token outputs.
Besides, there are also some earlier works noting asymmetries in the representation power and utility of different tokens during adaptation. For example, \citet{zhang2024dissecting} show that ``the adaptation of topic and style priors'' during finetuning are ``learned independently and primarily at the beginning of a text sequence.'' \citet{lin2024unlocking} also find that differences between aligned and unaligned base models introduced by the alignment fine-tuning vanishes as the sequence goes longer~(similar to the effect that we observe in Figure~\ref{fig:harmful-hexphi-kl}, though theirs are not in safety-specific contexts). Particularly, while \citet{lin2024unlocking} primarily tie such effects to question whether fine-tuning is even necessary to achieve current level of alignment~(e.g., in-context learning may already suffice to achieve a comparable level of alignment), we go much deeper into investigating the safety-specfic effects of this phenomenon and tie it to multiple downstream attacks and training-based mitigations. Besides, we find \citet{zhao2024weak} also note a similar token-wise effect, and they explicitly exploit this effect to design jailbreak attacks. Others have investigated fine-tuning dynamics through interpretability or pruning methods, which is distinct but somewhat related to the approach we take here~\citep{wei2024assessing,jain2023mechanistically}.

\textbf{Protecting The Safety Alignment at Initial Token Positions.} In Section~\ref{sec:constraining}, one important insight that motivates the design of our constrained fine-tuning loss is that ``safety alignment would be more difficult to be circumvented if we can protect the generative distribution of the model at the early token positions.'' We note that \citet{xu2024safedecoding} share a similar insight to ours in this regard. They find it is possible to design a defense against inference-time jailbreak attacks by simply identifying safety disclaimers and amplifying their token probabilities at the early token positions.

\textbf{Connections to Control Theory and Safe Reinforcement Learning.} Our data augmentation approach in Section~\ref{sec:make_alignment_deeper} relates to exploration requirements for optimal learning via policy gradient methods~\citep{agarwal2021theory}, learning recovery policies~\citep{thananjeyan2021recovery}, and safe control theory~\citep{burridge1999sequential,ames2019control}. We, however, leave deeper connections to this literature for future work.

\textbf{Other Notions of Safety Depth.} We also note that safety ``depth'' may be multi-dimensional in addition to token-based depth we describe here. Other considerations for depth would be the ability for models to retain safety properties after adaptation that some have previously discussed~\citep{henderson2023self,wei2024assessing,qi2023fine}.


\begin{comment}
A large body of work has examined improved methods for alignment~\citep{rafailov2023direct,ethayarajh2024kto,zou2023representation,bai2022constitutional,ouyang2022training,touvron2023llama-2,team2024gemma} and jailbreaking~\citep{andriushchenko2024jailbreaking,zou2023universal,qi2023fine,huang2023catastrophic}. Our work ties these potential failure modes of safety alignment methods to potential shortcuts taken during alignment optimization.
Some works have also noted asymmetries in the representation power and utility of different tokens during adaptation~\citep{zhang2024dissecting,lin2024unlocking,he2024s}. 
We build on this line of work to more deeply tie the dynamics of fine-tuning to downstream vulnerabilities and training-based defenses.
This is also an agenda-setting piece and we argue that alignment methods should optimize for deeper alignment.
See extended discussion in Appendix~\ref{appendix:related}.
\end{comment}

% For example, \citet{zhang2024dissecting} show that ``the adaptation of topic and style priors'' during finetuning are ``learned independently and primarily at the beginning of a text sequence.'' \citet{lin2024unlocking} also find that alignment fine-tuning primarily shifts the distribution of alignment only on the first few tokens, but primarily focus on in-context learning remedies. Our work builds on this line of work to more deeply tie the dynamics of fine-tuning to downstream vulnerabilities and training-based defenses. We take the position that alignment methods should seek to go beyond shallow alignment prevalent today. See extended discussion of related work in Appendix~\ref{appendix:related}.
% \citet{he2024s} also notes that to identify benign datapoints that break safety it was most useful to match gradients to harmful responses only for initial response tokens.

% Finally, our data augmentation approach relates to exploration requirements for optimal learning via policy gradient methods~\citep{agarwal2021theory}, learning recovery policies~\citep{thananjeyan2021recovery}, and safe control theory~\citep{burridge1999sequential,ames2019control}. We, however, leave deeper connections to this literature for future work.


%\xiangyu{I prefer deferring other alignment techniques to the related work section here. The "shallow alignment" we criticize is primarily just related to alignment fine-tuning~(e.g., SFT, RLHF), so I hesitate to mention decoding time intervention such as controlled decoding early on.} \ahmad{that's fair}

%\ahmad{might be good to also nudge to controlled decoding~\citep{yang2021fudge, mudgal2024controlled} or even best-of-$n$ at the end and ask whether such vulnerabilities carry over in those cases.}